\documentclass{../thesis}
% Font configuration for XeLaTeX/LuaLaTeX.
% Lao: Saysettha OT. English: Times New Roman.

\defaultfontfeatures{Ligatures=TeX,Scale=1.0}

\IfFontExistsTF{Saysettha OT}{
  \setmainfont[AutoFakeBold=2.0]{Saysettha OT}
  \newfontfamily\laofont[AutoFakeBold=2.0]{Saysettha OT}
}{
  \IfFontExistsTF{fonts/saysettha_ot.ttf}{
    \setmainfont[Path=fonts/,AutoFakeBold=2.0]{saysettha_ot.ttf}
    \newfontfamily\laofont[Path=fonts/,AutoFakeBold=2.0]{saysettha_ot.ttf}
  }{
    \IfFontExistsTF{../fonts/saysettha_ot.ttf}{
      \setmainfont[Path=../fonts/,AutoFakeBold=2.0]{saysettha_ot.ttf}
      \newfontfamily\laofont[Path=../fonts/,AutoFakeBold=2.0]{saysettha_ot.ttf}
    }{
      \IfFontExistsTF{Noto Sans Lao}{
        \setmainfont[AutoFakeBold=2.0]{Noto Sans Lao}
        \newfontfamily\laofont[AutoFakeBold=2.0]{Noto Sans Lao}
      }{
        \setmainfont{TeX Gyre Termes}
        \newfontfamily\laofont{TeX Gyre Termes}
      }
    }
  }
}

\IfFontExistsTF{Times New Roman}{
  \newfontfamily\englishfont{Times New Roman}
}{
  \newfontfamily\englishfont{TeX Gyre Termes}
}

\newcommand{\laotext}[1]{{\laofont #1}}
\newcommand{\engtext}[1]{{\englishfont #1}}


\begin{document}
\begin{titlepage}
\centering
\vspace*{3cm}
{\bfseries\fontsize{18pt}{24pt}\selectfont ຄູ່ມືການໃຊ້ LaTeX ສຳລັບນັກສຶກສາ\par}
\vspace{1cm}
{\engtext{\bfseries Lao Thesis Template User Guide}\par}
\vfill
{\fontsize{14pt}{18pt}\selectfont ສຳລັບແມ່ແບບບົດໂຄງການຈົບຊັ້ນ\par}
\end{titlepage}

\tableofcontents
\clearpage

\chapter{LaTeX ແມ່ນຫຍັງ}
LaTeX ແມ່ນລະບົບຈັດຮູບແບບເອກະສານທີ່ນິຍົມໃຊ້ໃນວຽກວິຊາການ ບົດວິຈັຍ ບົດຈົບ ແລະ ໜັງສື. ນັກສຶກສາຂຽນເນື້ອຫາໃນໄຟລ໌ \texttt{.tex} ແລ້ວໃຫ້ LaTeX ຈັດໜ້າ ເລກໜ້າ ສາລະບານ ຮູບ ຕາຕະລາງ ແລະ ບັນນານຸກົມໃຫ້ອັດຕະໂນມັດ.

\chapter{ການຕິດຕັ້ງ}
\section{TeX Live ຫຼື MiKTeX}
ຖ້າໃຊ້ Windows ສາມາດຕິດຕັ້ງ MiKTeX. ຖ້າໃຊ້ macOS ສາມາດຕິດຕັ້ງ MacTeX. ຖ້າໃຊ້ Linux ສາມາດຕິດຕັ້ງ TeX Live. ຫຼັງຈາກຕິດຕັ້ງແລ້ວ ກວດວ່າມີຄຳສັ່ງ \texttt{xelatex} ແລະ \texttt{biber}.

\section{Overleaf}
Overleaf ເປັນເວັບໄຊສຳລັບຂຽນ LaTeX ອອນລາຍ. ອັບໂຫຼດໂຟນເດີ \texttt{thesis-template} ເຂົ້າ Overleaf, ເປີດ Menu, ແລ້ວເລືອກ Compiler ເປັນ XeLaTeX.

\section{ຟອນ Saysettha OT}
ຕິດຕັ້ງ Saysettha OT ໃນລະບົບປະຕິບັດການ ຫຼື ວາງໄຟລ໌ຟອນໃນໂຟນເດີ \texttt{fonts/} ດ້ວຍຊື່ \texttt{saysettha\_ot.ttf}. ຟອນ Times New Roman ມັກຈະມີຢູ່ໃນ Windows; ຖ້າບໍ່ມີ ແມ່ແບບຈະໃຊ້ TeX Gyre Termes ແທນ.

\chapter{ການໃຊ້ແມ່ແບບ}
\section{ເປີດໄຟລ໌ main.tex}
ໄຟລ໌ຫຼັກຂອງບົດແມ່ນ \texttt{main.tex}. ໄຟລ໌ນີ້ຮວບຮວມປົກ ໜ້າຮັບຮອງ ບົດຄັດຫຍໍ້ ບົດທີ 1 ຫາ 5 ເອກະສານອ້າງອີງ ແລະ ພາກຜະນວກ.

\section{ແກ້ຂໍ້ມູນຫົວບົດ}
ເປີດ \texttt{config/settings.tex} ແລ້ວແກ້ຊື່ມະຫາວິທະຍາໄລ ຄະນະ ພາກວິຊາ ຊື່ບົດ ຊື່ນັກສຶກສາ ລະຫັດນັກສຶກສາ ອາຈານຜູ້ນຳພາ ແລະ ສົກສຶກສາ.

\section{ຂຽນບົດ}
ໄຟລ໌ບົດຢູ່ໃນໂຟນເດີ \texttt{chapters/}. ແກ້ໄຟລ໌ຕາມບົດທີ່ຕ້ອງການ. ໃຊ້ \verb|\chapter{...}| ສຳລັບບົດ, \verb|\section{...}| ສຳລັບຫົວຂໍ້ໃຫຍ່, ແລະ \verb|\subsection{...}| ສຳລັບຫົວຂໍ້ຍ່ອຍ.

\chapter{ຮູບ ແລະ ຕາຕະລາງ}
\section{ໃສ່ຮູບ}
ວາງຮູບໃນໂຟນເດີ \texttt{figures/} ແລ້ວໃຊ້:
\begin{verbatim}
\begin{figure}[H]
  \centering
  \includegraphics[width=0.7\textwidth]{figures/sample-figure.png}
  \caption{ຄຳອະທິບາຍຮູບ}
  \label{fig:myfigure}
\end{figure}
\end{verbatim}
ອ້າງຮູບດ້ວຍ \verb|\ref{fig:myfigure}|.

\section{ໃສ່ຕາຕະລາງ}
ໃຊ້ \texttt{table} ແລະ \texttt{tabularx} ສຳລັບຕາຕະລາງ. ຄຳອະທິບາຍຕາຕະລາງຈະຢູ່ເທິງຕາຕະລາງຕາມຮູບແບບແມ່ແບບ.

\chapter{ການອ້າງອີງ}
ເພີ່ມລາຍການອ້າງອີງໃນ \texttt{references/references.bib}. ໃນເນື້ອຫາໃຊ້ \verb|\cite{key}|. ຖ້າບັນນານຸກົມບໍ່ສະແດງ ໃຫ້ຄອມໄພລ໌ຕາມລຳດັບ XeLaTeX, Biber, XeLaTeX, XeLaTeX.

\chapter{ພາກຜະນວກ}
ໄຟລ໌ພາກຜະນວກຢູ່ໃນ \texttt{appendices/}. ຖ້າຕ້ອງການເພີ່ມພາກຜະນວກໃໝ່ ໃຫ້ສ້າງໄຟລ໌ \texttt{.tex} ແລ້ວເພີ່ມ \verb|\input{appendices/filename}| ໃນ \texttt{main.tex}.

\chapter{ການສ້າງ PDF}
ຖ້າໃຊ້ command line ໃຫ້ພິມ:
\begin{verbatim}
make pdf
\end{verbatim}
ຖ້າບໍ່ມີ Make ໃຫ້ຮັນ:
\begin{verbatim}
xelatex main.tex
biber main
xelatex main.tex
xelatex main.tex
\end{verbatim}

\chapter{ບັນຫາທີ່ພົບເລື້ອຍ}
\section{Font not found}
ຖ້າພົບຂໍ້ຜິດພາດຟອນ ໃຫ້ຕິດຕັ້ງ Saysettha OT ຫຼື ວາງໄຟລ໌ \texttt{saysettha\_ot.ttf} ໃນ \texttt{fonts/}.

\section{ພາສາລາວສະແດງບໍ່ຖືກ}
ກວດວ່າໃຊ້ XeLaTeX ຫຼື LuaLaTeX ບໍ່ແມ່ນ pdfLaTeX. ກວດວ່າໄຟລ໌ບັນທຶກເປັນ UTF-8.

\section{Bibliography ບໍ່ສະແດງ}
ຮັນ \texttt{biber main} ຫຼັງຈາກຮັນ XeLaTeX ຄັ້ງທຳອິດ ແລ້ວຮັນ XeLaTeX ອີກສອງຄັ້ງ.

\section{Image not found}
ກວດຊື່ໄຟລ໌ ນາມສະກຸນໄຟລ໌ ແລະ path ໃຫ້ຖືກຕ້ອງ. ແນະນຳໃຫ້ໃຊ້ຊື່ໄຟລ໌ພາສາອັງກິດບໍ່ມີຊ່ອງວ່າງ.

\chapter{ກ່ອນສົ່ງບົດ}
\begin{enumerate}
  \item ກວດຊື່ບົດ ຊື່ນັກສຶກສາ ແລະ ອາຈານຜູ້ນຳພາ.
  \item ກວດສາລະບານ ສາລະບານຮູບ ແລະ ສາລະບານຕາຕະລາງ.
  \item ກວດການອ້າງອີງທຸກລາຍການ.
  \item ກວດວ່າ PDF ສະແດງພາສາລາວຖືກຕ້ອງ.
  \item ສົ່ງ PDF ສຸດທ້າຍຕາມຂໍ້ກຳນົດຂອງພາກວິຊາ.
\end{enumerate}

\end{document}
